\import{function.tex}

\section{Rings and fields}

The ring hierarchy uses operation names distinct from the generic multiplicative
group interface.  This keeps additive and multiplicative structure available
simultaneously without overloading the single group operation field.

\begin{struct}\label{ring}
    A ring $R$ is a onesorted structure equipped with
    \begin{enumerate}
        \item\label{ring_add_graph}
            an addition
            $\ringAdd[R]:\carrier[R]\times\carrier[R]\to\carrier[R]$;
        \item\label{ring_zero_type}
            a zero element $\ringZero[R]\in\carrier[R]$;
        \item\label{ring_neg_graph}
            an additive inverse
            $\ringNeg[R]:\carrier[R]\to\carrier[R]$;
        \item\label{ring_mul_graph}
            a multiplication
            $\ringMul[R]:\carrier[R]\times\carrier[R]\to\carrier[R]$;
        \item\label{ring_one_type}
            a unit element $\ringOne[R]\in\carrier[R]$;
    \end{enumerate}
    such that
    \begin{enumerate}
        \item\label{ring_add_associative}
            for all $x,y,z\in\carrier[R]$ we have
                $\ringAdd[R](x,\ringAdd[R](y,z))
                    =\ringAdd[R](\ringAdd[R](x,y),z)$.
        \item\label{ring_add_commutative}
            for all $x,y\in\carrier[R]$ we have
                $\ringAdd[R](x,y)=\ringAdd[R](y,x)$.
        \item\label{ring_add_zero_left}
            for all $x\in\carrier[R]$ we have
                $\ringAdd[R](\ringZero[R],x)=x$.
        \item\label{ring_add_inverse_left}
            for all $x\in\carrier[R]$ we have
                $\ringAdd[R](\ringNeg[R](x),x)=\ringZero[R]$.
        \item\label{ring_mul_associative}
            for all $x,y,z\in\carrier[R]$ we have
                $\ringMul[R](x,\ringMul[R](y,z))
                    =\ringMul[R](\ringMul[R](x,y),z)$.
        \item\label{ring_mul_one_left}
            for all $x\in\carrier[R]$ we have
                $\ringMul[R](\ringOne[R],x)=x$.
        \item\label{ring_mul_one_right}
            for all $x\in\carrier[R]$ we have
                $\ringMul[R](x,\ringOne[R])=x$.
        \item\label{ring_distributive_left}
            for all $x,y,z\in\carrier[R]$ we have
                $\ringMul[R](x,\ringAdd[R](y,z))
                    =\ringAdd[R]
                        (\ringMul[R](x,y),\ringMul[R](x,z))$.
        \item\label{ring_distributive_right}
            for all $x,y,z\in\carrier[R]$ we have
                $\ringMul[R](\ringAdd[R](x,y),z)
                    =\ringAdd[R]
                        (\ringMul[R](x,z),\ringMul[R](y,z))$.
    \end{enumerate}
\end{struct}

\begin{theorem}\label{exact_ring_extensionality}
    Let $R,S$ be sets.
    Suppose $R$ is an exact ring and $S$ is an exact ring.
    Suppose $\carrier[R]=\carrier[S]$.
    Suppose $\ringAdd[R]=\ringAdd[S]$.
    Suppose $\ringZero[R]=\ringZero[S]$.
    Suppose $\ringNeg[R]=\ringNeg[S]$.
    Suppose $\ringMul[R]=\ringMul[S]$.
    Suppose $\ringOne[R]=\ringOne[S]$.
    Then $R=S$.
\end{theorem}
\begin{proof}
    Follows by structure extensionality.
\end{proof}

\begin{theorem}\label{exact_ring_equality_from_fields}
    Let $R,S$ be sets.
    Suppose $R$ is an exact ring and $S$ is an exact ring.
    Suppose $\carrier[R]=\carrier[S]$.
    Suppose $\ringAdd[R]=\ringAdd[S]$.
    Suppose $\ringZero[R]=\ringZero[S]$.
    Suppose $\ringNeg[R]=\ringNeg[S]$.
    Suppose $\ringMul[R]=\ringMul[S]$.
    Suppose $\ringOne[R]=\ringOne[S]$.
    Then $R=S$.
\end{theorem}
\begin{proof}
    Follows by \obv{exact_ring_extensionality}.
\end{proof}

The explicit structure index is part of the notation.  It keeps expressions
unambiguous when several ring structures occur in the same argument.

\begin{abbreviation}\label{ring_add_notation}
    %! infixl 1
    $x+_{R}y=\ringAdd[R](x,y)$.
\end{abbreviation}

\begin{abbreviation}\label{ring_zero_notation}
    $0_{R}=\ringZero[R]$.
\end{abbreviation}

\begin{abbreviation}\label{ring_one_notation}
    $1_{R}=\ringOne[R]$.
\end{abbreviation}

\begin{abbreviation}\label{ring_negation_notation}
    $-_{R}x=\ringNeg[R](x)$.
\end{abbreviation}

\begin{proposition}\label{ring_add_type}
    Let $R$ be a ring.
    Let $x,y\in\carrier[R]$.
    Then $\ringAdd[R](x,y)\in\carrier[R]$.
\end{proposition}
\begin{proof}
    $(x,y)\in\carrier[R]\times\carrier[R]$
        by \obv{times_tuple_intro}.
    Follows by \obv{ring_add_graph,funs_type_apply}.
\end{proof}

\begin{proposition}\label{ring_neg_type}
    Let $R$ be a ring.
    Let $x\in\carrier[R]$.
    Then $\ringNeg[R](x)\in\carrier[R]$.
\end{proposition}
\begin{proof}
    Follows by \obv{ring_neg_graph,funs_type_apply}.
\end{proof}

\begin{proposition}\label{ring_mul_type}
    Let $R$ be a ring.
    Let $x,y\in\carrier[R]$.
    Then $\ringMul[R](x,y)\in\carrier[R]$.
\end{proposition}
\begin{proof}
    $(x,y)\in\carrier[R]\times\carrier[R]$
        by \obv{times_tuple_intro}.
    Follows by \obv{ring_mul_graph,funs_type_apply}.
\end{proof}

\begin{struct}\label{commutative_ring}
    A commutative ring $R$ is a ring such that
    \begin{enumerate}
        \item\label{commutative_ring_mul_commutative}
            for all $x,y\in\carrier[R]$ we have
                $\ringMul[R](x,y)=\ringMul[R](y,x)$.
    \end{enumerate}
\end{struct}

\begin{proposition}\label{commutative_ring_intro}
    Let $R$ be a ring.
    Suppose for all $x,y\in\carrier[R]$ we have
        $\ringMul[R](x,y)=\ringMul[R](y,x)$.
    Then $R$ is a commutative ring.
\end{proposition}
\begin{proof}
    Follows by definition.
\end{proof}

\begin{struct}\label{integral_domain}
    An integral domain $R$ is a commutative ring such that
    \begin{enumerate}
        \item\label{integral_domain_nontrivial}
            $\ringZero[R]\neq\ringOne[R]$.
        \item\label{integral_domain_no_zero_divisors}
            for all $x,y\in\carrier[R]$ such that
                $\ringMul[R](x,y)=\ringZero[R]$ we have
                $x=\ringZero[R]$ or $y=\ringZero[R]$.
    \end{enumerate}
\end{struct}

\begin{proposition}\label{integral_domain_intro}
    Let $R$ be a commutative ring.
    Suppose $\ringZero[R]\neq\ringOne[R]$.
    Suppose for all $x,y\in\carrier[R]$ such that
        $\ringMul[R](x,y)=\ringZero[R]$ we have
        $x=\ringZero[R]$ or $y=\ringZero[R]$.
    Then $R$ is an integral domain.
\end{proposition}
\begin{proof}
    Follows by definition.
\end{proof}

\begin{struct}\label{division_ring}
    A division ring $R$ is a ring equipped with
    \begin{enumerate}
        \item\label{division_ring_inverse_graph}
            a multiplicative inverse
            $\ringInv[R]:\carrier[R]\to\carrier[R]$;
    \end{enumerate}
    such that
    \begin{enumerate}
        \item\label{division_ring_nontrivial}
            $\ringZero[R]\neq\ringOne[R]$.
        \item\label{division_ring_inverse_zero}
            $\ringInv[R](\ringZero[R])=\ringZero[R]$.
        \item\label{division_ring_inverse_left}
            for all $x\in\carrier[R]$ such that
                $x\neq\ringZero[R]$ we have
                $\ringMul[R](\ringInv[R](x),x)=\ringOne[R]$.
        \item\label{division_ring_inverse_right}
            for all $x\in\carrier[R]$ such that
                $x\neq\ringZero[R]$ we have
                $\ringMul[R](x,\ringInv[R](x))=\ringOne[R]$.
    \end{enumerate}
\end{struct}

Division-ring inversion is a total operation on the carrier.  Its value at
zero is fixed to zero, while the left and right inverse laws are deliberately
guarded by nonzeroness.

\begin{proposition}\label{division_ring_inverse_type}
    Let $R$ be a division ring.
    Let $x\in\carrier[R]$.
    Then $\ringInv[R](x)\in\carrier[R]$.
\end{proposition}
\begin{proof}
    Follows by \obv{division_ring_inverse_graph,funs_type_apply}.
\end{proof}

\begin{struct}\label{field}
    A field $F$ is a division ring such that
    \begin{enumerate}
        \item\label{field_mul_commutative}
            for all $x,y\in\carrier[F]$ we have
                $\ringMul[F](x,y)=\ringMul[F](y,x)$.
    \end{enumerate}
\end{struct}

\begin{corollary}\label{commutative_ring_is_ring}
    Let $R$ be a commutative ring.
    Then $R$ is a ring.
\end{corollary}
\begin{proof}
    Follows by \obv{commutative_ring}.
\end{proof}

\begin{corollary}\label{integral_domain_is_commutative_ring}
    Let $R$ be an integral domain.
    Then $R$ is a commutative ring.
\end{corollary}
\begin{proof}
    Follows by \obv{integral_domain}.
\end{proof}

\begin{corollary}\label{division_ring_is_ring}
    Let $R$ be a division ring.
    Then $R$ is a ring.
\end{corollary}
\begin{proof}
    Follows by \obv{division_ring}.
\end{proof}

\begin{corollary}\label{field_is_division_ring}
    Let $F$ be a field.
    Then $F$ is a division ring.
\end{corollary}
\begin{proof}
    Follows by \obv{field}.
\end{proof}

\begin{proposition}\label{ring_add_zero_right}
    Let $R$ be a ring.
    Let $x\in\carrier[R]$.
    Then $\ringAdd[R](x,\ringZero[R])=x$.
\end{proposition}
\begin{proof}
    Follows by
        \auto{ring_zero_type,ring_add_commutative,ring_add_zero_left}.
\end{proof}

\begin{proposition}\label{ring_add_inverse_right}
    Let $R$ be a ring.
    Let $x\in\carrier[R]$.
    Then $\ringAdd[R](x,\ringNeg[R](x))=\ringZero[R]$.
\end{proposition}
\begin{proof}
    Follows by
        \auto{ring_neg_type,ring_add_commutative,ring_add_inverse_left}.
\end{proof}

\begin{theorem}\label{ring_add_cancel_left}
    Let $R$ be a ring.
    Let $a,b,c\in\carrier[R]$.
    Suppose $\ringAdd[R](a,b)=\ringAdd[R](a,c)$.
    Then $b=c$.
\end{theorem}
\begin{proof}
    $\ringNeg[R](a)\in\carrier[R]$ by \obv{ring_neg_type}.
    $\ringAdd[R](\ringNeg[R](a),\ringAdd[R](a,b))
        =\ringAdd[R](\ringNeg[R](a),\ringAdd[R](a,c))$.
    Hence
    $\ringAdd[R](\ringAdd[R](\ringNeg[R](a),a),b)
        =\ringAdd[R](\ringAdd[R](\ringNeg[R](a),a),c)$
        by \obv{ring_add_associative,ring_add_type}.
    Follows by
        \auto{ring_add_inverse_left,ring_add_zero_left}.
\end{proof}

\begin{theorem}\label{ring_add_cancel_right}
    Let $R$ be a ring.
    Let $a,b,c\in\carrier[R]$.
    Suppose $\ringAdd[R](a,c)=\ringAdd[R](b,c)$.
    Then $a=b$.
\end{theorem}
\begin{proof}
    Follows by
        \auto{ring_add_commutative,ring_add_cancel_left}.
\end{proof}

\begin{theorem}\label{ring_neg_unique}
    Let $R$ be a ring.
    Let $x,y\in\carrier[R]$.
    Suppose $\ringAdd[R](x,y)=\ringZero[R]$.
    Then $y=\ringNeg[R](x)$.
\end{theorem}
\begin{proof}
    $\ringNeg[R](x)\in\carrier[R]$ by \obv{ring_neg_type}.
    $\ringAdd[R](x,\ringNeg[R](x))=\ringZero[R]$
        by \obv{ring_add_inverse_right}.
    Follows by \obv{ring_add_cancel_left}.
\end{proof}

\begin{proposition}\label{ring_neg_zero}
    Let $R$ be a ring.
    Then $\ringNeg[R](\ringZero[R])=\ringZero[R]$.
\end{proposition}
\begin{proof}
    $\ringZero[R]\in\carrier[R]$ by \obv{ring_zero_type}.
    $\ringNeg[R](\ringZero[R])\in\carrier[R]$
        by \obv{ring_neg_type}.
    $\ringAdd[R](\ringZero[R],\ringNeg[R](\ringZero[R]))
        =\ringNeg[R](\ringZero[R])$
        by \obv{ring_add_zero_left}.
    $\ringAdd[R](\ringZero[R],\ringNeg[R](\ringZero[R]))
        =\ringZero[R]$
        by \obv{ring_add_inverse_right}.
\end{proof}

\begin{theorem}\label{ring_neg_involutive}
    Let $R$ be a ring.
    Let $x\in\carrier[R]$.
    Then $\ringNeg[R](\ringNeg[R](x))=x$.
\end{theorem}
\begin{proof}
    $\ringNeg[R](x)\in\carrier[R]$ by \obv{ring_neg_type}.
    $\ringAdd[R](\ringNeg[R](x),x)=\ringZero[R]$
        by \obv{ring_add_inverse_left}.
    Follows by \obv{ring_neg_unique}.
\end{proof}

\begin{theorem}\label{ring_mul_zero_right}
    Let $R$ be a ring.
    Let $x\in\carrier[R]$.
    Then $\ringMul[R](x,\ringZero[R])=\ringZero[R]$.
\end{theorem}
\begin{proof}
    Let $a=\ringMul[R](x,\ringZero[R])$.
    $\ringZero[R]\in\carrier[R]$ by \obv{ring_zero_type}.
    $a\in\carrier[R]$ by \auto{ring_zero_type,ring_mul_type}.
    $\ringAdd[R](\ringZero[R],\ringZero[R])=\ringZero[R]$
        by \obv{ring_add_zero_left,ring_zero_type}.
    Hence $a=\ringMul[R]
        (x,\ringAdd[R](\ringZero[R],\ringZero[R]))$.
    Hence $a=\ringAdd[R](a,a)$
        by \auto{ring_distributive_left}.
    $\ringAdd[R](a,\ringZero[R])=a$
        by \obv{ring_add_zero_right}.
    Hence
        $\ringAdd[R](a,\ringZero[R])=\ringAdd[R](a,a)$.
    Hence $\ringZero[R]=a$ by \obv{ring_add_cancel_left}.
\end{proof}

\begin{theorem}\label{ring_mul_zero_left}
    Let $R$ be a ring.
    Let $x\in\carrier[R]$.
    Then $\ringMul[R](\ringZero[R],x)=\ringZero[R]$.
\end{theorem}
\begin{proof}
    Let $a=\ringMul[R](\ringZero[R],x)$.
    $\ringZero[R]\in\carrier[R]$ by \obv{ring_zero_type}.
    $a\in\carrier[R]$ by \auto{ring_zero_type,ring_mul_type}.
    $\ringAdd[R](\ringZero[R],\ringZero[R])=\ringZero[R]$
        by \obv{ring_add_zero_left,ring_zero_type}.
    Hence $a=\ringMul[R]
        (\ringAdd[R](\ringZero[R],\ringZero[R]),x)$.
    Hence $a=\ringAdd[R](a,a)$
        by \auto{ring_distributive_right}.
    $\ringAdd[R](a,\ringZero[R])=a$
        by \obv{ring_add_zero_right}.
    Hence
        $\ringAdd[R](a,\ringZero[R])=\ringAdd[R](a,a)$.
    Hence $\ringZero[R]=a$ by \obv{ring_add_cancel_left}.
\end{proof}

\begin{theorem}\label{ring_neg_product_left}
    Let $R$ be a ring.
    Let $x,y\in\carrier[R]$.
    Then
        $\ringMul[R](\ringNeg[R](x),y)
            =\ringNeg[R](\ringMul[R](x,y))$.
\end{theorem}
\begin{proof}
    $\ringNeg[R](x)\in\carrier[R]$ by \obv{ring_neg_type}.
    $\ringMul[R](x,y)$ and
        $\ringMul[R](\ringNeg[R](x),y)$ are elements of $\carrier[R]$
        by \obv{ring_mul_type}.
    $\ringAdd[R](\ringNeg[R](x),x)=\ringZero[R]$
        by \obv{ring_add_inverse_left}.
    Hence
    $\ringAdd[R]
        (\ringMul[R](\ringNeg[R](x),y),\ringMul[R](x,y))
        =\ringZero[R]$
        by \auto{ring_distributive_right,ring_mul_zero_left}.
    Hence
    $\ringAdd[R]
        (\ringMul[R](x,y),\ringMul[R](\ringNeg[R](x),y))
        =\ringZero[R]$
        by \obv{ring_add_commutative}.
    Follows by \obv{ring_neg_unique}.
\end{proof}

\begin{theorem}\label{ring_neg_product_right}
    Let $R$ be a ring.
    Let $x,y\in\carrier[R]$.
    Then
        $\ringMul[R](x,\ringNeg[R](y))
            =\ringNeg[R](\ringMul[R](x,y))$.
\end{theorem}
\begin{proof}
    $\ringNeg[R](y)\in\carrier[R]$ by \obv{ring_neg_type}.
    $\ringMul[R](x,y)$ and
        $\ringMul[R](x,\ringNeg[R](y))$ are elements of $\carrier[R]$
        by \obv{ring_mul_type}.
    $\ringAdd[R](\ringNeg[R](y),y)=\ringZero[R]$
        by \obv{ring_add_inverse_left}.
    Hence
    $\ringAdd[R]
        (\ringMul[R](x,\ringNeg[R](y)),\ringMul[R](x,y))
        =\ringZero[R]$
        by \auto{ring_distributive_left,ring_mul_zero_right}.
    Hence
    $\ringAdd[R]
        (\ringMul[R](x,y),\ringMul[R](x,\ringNeg[R](y)))
        =\ringZero[R]$
        by \obv{ring_add_commutative}.
    Follows by \obv{ring_neg_unique}.
\end{proof}

\begin{theorem}\label{division_ring_inverse_unique}
    Let $R$ be a division ring.
    Let $x,y\in\carrier[R]$.
    Suppose $x\neq\ringZero[R]$ and
        $\ringMul[R](x,y)=\ringOne[R]$.
    Then $y=\ringInv[R](x)$.
\end{theorem}
\begin{proof}
    $R$ is a ring by \obv{division_ring_is_ring}.
    $\ringInv[R](x)\in\carrier[R]$
        by \obv{division_ring_inverse_type}.
    $\ringMul[R](\ringInv[R](x),
        \ringMul[R](x,y))
        =\ringMul[R](\ringInv[R](x),\ringOne[R])$.
    Hence
    $\ringMul[R](
        \ringMul[R](\ringInv[R](x),x),y)
        =\ringInv[R](x)$
        by \obv{ring_mul_associative,ring_mul_one_right,ring_mul_type}.
    Follows by
        \auto{division_ring_inverse_left,ring_mul_one_left}.
\end{proof}

\begin{proposition}\label{division_ring_inverse_one}
    Let $R$ be a division ring.
    Then $\ringInv[R](\ringOne[R])=\ringOne[R]$.
\end{proposition}
\begin{proof}
    $\ringOne[R]\in\carrier[R]$ by \obv{ring_one_type}.
    $\ringOne[R]\neq\ringZero[R]$
        by \obv{division_ring_nontrivial}.
    $\ringMul[R](\ringOne[R],\ringOne[R])=\ringOne[R]$
        by \obv{ring_mul_one_left}.
    Follows by \obv{division_ring_inverse_unique}.
\end{proof}

\begin{theorem}\label{division_ring_no_zero_divisors}
    Let $R$ be a division ring.
    Let $x,y\in\carrier[R]$.
    Suppose $\ringMul[R](x,y)=\ringZero[R]$.
    Then $x=\ringZero[R]$ or $y=\ringZero[R]$.
\end{theorem}
\begin{proof}
    $R$ is a ring by \obv{division_ring_is_ring}.
    \begin{byCase}
        \caseOf{$x=\ringZero[R]$.}
            Follows.
        \caseOf{$x\neq\ringZero[R]$.}
            $\ringInv[R](x)\in\carrier[R]$
                by \obv{division_ring_inverse_type}.
            $\ringMul[R](\ringInv[R](x),\ringMul[R](x,y))
                =\ringMul[R](\ringInv[R](x),\ringZero[R])$.
            Hence
            $\ringMul[R](
                \ringMul[R](\ringInv[R](x),x),y)
                =\ringZero[R]$
                by
                \obv{ring_mul_associative,ring_mul_zero_right,ring_mul_type}.
            Hence $y=\ringZero[R]$
                by \auto{division_ring_inverse_left,ring_mul_one_left}.
    \end{byCase}
\end{proof}

\begin{theorem}\label{field_is_commutative_ring}
    Let $F$ be a field.
    Then $F$ is a commutative ring.
\end{theorem}
\begin{proof}
    $F$ is a division ring by \obv{field_is_division_ring}.
    Hence $F$ is a ring by \obv{division_ring_is_ring}.
    Follows by \obv{commutative_ring_intro,field_mul_commutative}.
\end{proof}

\begin{theorem}\label{field_is_integral_domain}
    Let $F$ be a field.
    Then $F$ is an integral domain.
\end{theorem}
\begin{proof}
    $F$ is a division ring by \obv{field_is_division_ring}.
    $F$ is a commutative ring by \obv{field_is_commutative_ring}.
    $\ringZero[F]\neq\ringOne[F]$
        by \obv{division_ring_nontrivial}.
    Follows by
        \obv{integral_domain_intro,division_ring_no_zero_divisors}.
\end{proof}
